\section{My impact}

\textbf{3a. How my contributions drove the team forward.} The honest thesis of this section is the one I wrote into our group complaint letter: \emph{the simulation framework we built ourselves is the only reason this project produced any verifiable engineering output at all}. I want to defend that claim with mechanism, not assertion, because it is exactly the kind of line that at Merit-level reads as self-flattery.

The anchor is a full Kolb (1984) cycle, which I used deliberately rather than retrospectively. My \emph{concrete experience} was three collapsed flight days: the 2026-03-11 weather cancellation, the 2026-03-30 hardware session that produced no airborne minutes, and a third day that ended with our pilot standing on a wet field while Edward and I debugged a corrupt \texttt{calibration\_data.npz}. After the first cancellation my \emph{reflective observation} was harder than ``we lost a day'': the team's entire definition of ``progress'' had silently assumed hardware access we did not have, and the rest of the team was idle because their planned work needed a drone in the air. I was blocked too, but I had a laptop. That asymmetry was the data point. My \emph{abstract conceptualisation} was the move I am most proud of: if \texttt{vision.py} exposed a single interface (\texttt{detect\_in\_image} returning \texttt{(found, x, y, conf)}) and auto-selected between Ultralytics on laptop and TFLite on Pi, and if \texttt{config.py} auto-detected its hardware, then the same state machine, geofence, and search pattern could run in both places --- anything that ran in simulation would, in principle, run on hardware. My \emph{active experimentation} was six weeks of hammering on that premise: \texttt{simple\_simulator.py} grew to 2508 lines with GPS drift and camera shake; \texttt{main.py} was refactored 1411\,$\to$\,754 lines across six modules with 146 pytest tests; and the \texttt{tests/flight/0a--4} ladder plus 127 unit tests gave me a way to know, before any field day, which scripts could survive contact with hardware. When we finally ran on the Pi the delta was thin --- a bbox-coordinate bug (commit \texttt{3d62e6a}), a double-scaling fix (\texttt{7eb7cc8}) and a deleted calibration file --- and the system came up. I now see the simulation framework was not a backup plan but a \emph{statement}: verifiable engineering is possible without hardware access, and simulation-first is no longer a personal preference for me but a claim about what the discipline owes.

Three other contributions matter here and I mention them compactly. I unblocked Robin by building him a dedicated HTTP integration path (\texttt{passive\_watch\_2.py}, commit \texttt{a5b1ab7}) with a JSON schema and \texttt{--smart-clear} flag. I unblocked Demetro by building his two FDR slides from scratch and writing 465 words of speaking script across 4--5 iterations. And I drove the group's external communication about the hardware failures via the six-iteration complaint letter the whole team signed.

\emph{Methods evaluated against observable behaviour change per author-hour.} I used six comms methods with named audiences during this project: (1) a 4-hour \texttt{VISION\_PIPELINE\_FOR\_COLLEAGUE.md} doc for Robin; (2) the \texttt{--robin} JSON/HTTP interface for Robin's mission script; (3) FDR slides and speaking scripts for Demetro (and the FDR panel); (4) the browser ground-station for the pilot; (5) the clarification letter for the course organiser; and (6) a 400-word Teams pivot message for the team. Against the criterion \emph{did a teammate change behaviour within 48 hours}, the two fastest were the pivot message (team committed to simulation-first the same night) and the \texttt{--robin} flag (Robin integrated unaided for two days). The 3000-word handoff doc, by contrast, sat unread until Robin finally needed it in wk20. The lesson, which I did not expect, is that prose quality was almost orthogonal to behaviour change --- timing, pairing, and interface shape dominated.

But the sober closing of 3a is that the simulation framework could \emph{not} tell us what it could not tell us. It could not predict motion blur at 5\,m/s with a real IMX296. It could not find the descent stall we accepted as a SITL artefact. It could not tell us what GPS timing lag would do to a multi-pass lock --- that came out of DJI video analysis, after the fact. The statement my simulation made is true, and the limits of that statement are also true: sim-to-real parity is an argument I can only win about the things I modelled, and I was disciplined about that only because I had been forced to be.

\textbf{3b. What I would add or re-focus on in future.} I mistook technical leadership for team leadership. I now see they are orthogonal --- you can build great tools that do not enable anyone --- and in this project the gap showed up as a metric: the number of times a teammate other than me edited \texttt{simple\_simulator.py} or \texttt{main.py} was zero. Three concrete mechanisms, each with a falsifiability test, follow. First, I will shift from ``building tools I need'' to ``building tools the team needs,'' and each new tool must have at least one named teammate beneficiary documented \emph{before} it ships --- if I cannot name them, I do not start. Second, no file larger than 500 lines will be refactored solo: I will pair on the review, and the observable signal is whether the pair partner can edit that file independently within 48 hours of the walkthrough. Third, I will schedule two hours per week of explicit teaching sessions, graded not by whether I explained something but by whether the teammate completes a handoff task within 48 hours without asking me a second question. If after four weeks the signals are not firing, the mechanism must change, not my reassurance that I meant well.
